\chapter{Background}
\label{chap:background}

\section{Monte Carlo Option Pricing}

Monte Carlo simulation estimates an expectation by averaging random samples. In option pricing, this usually means simulating asset paths under a risk-neutral model and discounting the payoff. For a European call under geometric Brownian motion (GBM), the terminal asset value can be sampled directly and compared against the Black--Scholes analytical price \cite{black1973pricing}. More complex payoffs, such as Asian options, require path simulation because the payoff depends on the realised path rather than only the terminal value. Monte Carlo methods have long been used for option pricing because they are flexible and naturally parallel across paths \cite{boyle1977options}.

The flexibility comes at a cost. Accuracy improves slowly with the number of paths: the standard error decreases at the usual \(O(M^{-1/2})\) Monte Carlo rate, where \(M\) is the path count. A benchmark therefore needs to distinguish numerical error from implementation overhead. It must also ensure that different engines are pricing comparable workloads with comparable random seeds, time grids and model parameters.

\section{Automatic Differentiation}

Automatic differentiation evaluates derivatives by applying the chain rule to the operations performed by a program \cite{griewank2008evaluating,baydin2018automatic}. It differs from finite differences because it does not require manually perturbing inputs and rerunning the whole computation for each sensitivity. In this project, AD is used to compute Greeks for JAX implementations of the Monte Carlo workloads.

Forward-mode AD propagates tangent information from inputs to outputs. It is attractive when the number of differentiated inputs is small. Reverse-mode AD propagates adjoints from outputs back to inputs. It is often efficient for scalar outputs with many inputs, but it can require retaining or reconstructing intermediate values. In simulation workloads, this trade-off interacts with path count, time steps, vectorisation and JIT compilation. Measuring the overhead empirically is therefore more reliable than assuming that one AD mode is always best.

\section{Engines and Runtime Effects}

The evaluated engines represent different implementation styles:
\begin{itemize}
    \item \textbf{NumPy CPU} provides a straightforward vectorised Python baseline.
    \item \textbf{JAX/XLA} provides JIT compilation and AD transformations.
    \item \textbf{C++/OpenMP} represents an ahead-of-time compiled native extension with thread-level parallelism.
    \item \textbf{Rust/Rayon} represents a memory-safe compiled implementation with data-parallel execution.
\end{itemize}

These engines differ in compilation overhead, memory layout, random-number handling, thread scheduling and support for AD. They also differ by workload. A direct European payoff can expose different bottlenecks from a local-volatility workload with many time steps, and an Asian option introduces path dependence. This motivates a framework-level comparison rather than isolated microbenchmarks.

\section{Benchmarking and Reproducibility}

A useful benchmark should be repeatable, inspectable and fair. Repeatability requires fixed configuration, deterministic seeding where possible and consistent warm-up. Inspectability requires storing not only aggregate results but also workload parameters, engine identifiers, AD modes, environment metadata and code version. Fairness requires comparing like with like: the same payoff, path count and model parameters should be used across engines wherever those engines support the workload.

Cloud benchmarking adds further complications. Runtime is affected by machine type, region, zone, library builds and background variability. Cost-performance also depends on instance pricing. For this reason, this project stores cloud metadata alongside timings and treats cloud results as measured evidence from specific instances rather than universal statements about all hardware.

\section{Successive Halving for Benchmark Selection}

Successive Halving is a budget allocation strategy originally developed for selecting promising configurations by evaluating many candidates cheaply, discarding weaker ones, and spending more budget on the survivors \cite{jamieson2016nonstochastic}. Hyperband generalises this idea across multiple brackets of resource allocation \cite{li2016hyperband}. The key idea is simple: if cheap low-budget measurements are informative, they can reduce the number of expensive full-budget evaluations.

This project adapts that idea to cloud benchmarking. The low-budget resource is the Monte Carlo path count used for probe runs, and the expensive resource is a full benchmark run at the target path count. The SHA profiler ranks configurations from cheap probes, promotes promising candidates and then compares the final selections against a full-grid baseline using Pareto recovery, runtime regret, cost regret and rank correlation. The adaptation is empirical: the report evaluates whether the probes were predictive for the measured workloads rather than assuming they must be.
